\begin{flushleft}
\renewcommand{\arraystretch}{1.3}
\begin{tabular}{@{}p{0.32\linewidth}p{0.62\linewidth}@{}}
\hline
Thesis topic:                                                 & \textbf{Improving situational awareness of autonomous vessels using computer vision} \\
Thesis topic in Estonian:                                     & \textbf{Autonoomse laeva situatsiooniteadlikkuse arendamine masinnägemise abil} \\
Student:                                                      & \textbf{Tanel Treuberg, 193761EAAB} \\
Specialty:                                                    & \textbf{Electrical Power Engineering and Mechatronics} \\
Type of thesis:                                               & \textbf{Bachelor's thesis} \\
Thesis supervisor:                                            & \textbf{Heigo Mõlder, PhD} \\
Thesis co-supervisor:                                         & \textbf{Karl Janson, PhD} \\
(organisation, position, contact)                             & \textbf{TalTech Department of Computer Systems, e-mail: karl.janson@taltech.ee} \\
Assignment validity period (set by the supervisor):           & 2022/2023 Spring \\
Thesis submission deadline:                                   & \textbf{18.05.23} \\
\hline
\end{tabular}
\end{flushleft}

\vspace{1.5em}

\noindent\rule{6cm}{0.4pt}\\
Student (signature)

\noindent\rule{6cm}{0.4pt}\\
Supervisor (signature)

\noindent\rule{6cm}{0.4pt}\\
Programme director (signature)

\noindent\rule{6cm}{0.4pt}\\
Co-supervisor (signature)

\section*{1. Rationale}

Within a joint project between the TalTech Department of Electrical Power Engineering and Mechatronics, the TalTech Centre of Excellence in Small-Craft Building, and the maritime company Baltic Workboats, the aim is to automate vessel sea trials that have until now been carried out by various captains with manual control. The results of the trials are therefore not objective; they depend on how each captain handles the vessel, and deviations can appear in the data on whose basis it is not possible to confirm the vessel's parameters with the desired accuracy. Conducting the trials is also very time-consuming, since a particular manoeuvre often has to be repeated several times because in some runs the vessel is steered with too large a radius or at the wrong speed -- human error introduces unwelcome deviations in the trial results.

Within the cooperation project, these sea trials are planned to be performed with the help of a computer program, in which the vessel's situational awareness plays a key role. This thesis aims to develop the vessel's situational awareness by building a computer-vision-based solution for perceiving the surroundings.

\section*{2. Aim of the work}

The aim of the thesis is to develop a computer-vision system suitable for the automated conduct of vessel sea trials. This means selecting suitable hardware and creating the software for the computer-vision system, with intermediate testing in the harbour and finally on a real vessel. It also includes investigating how the software can be optimised so as to make the most efficient use of the hardware, and as a result make the software's processing smoother and faster.

\section*{3. List of questions to be addressed}

\begin{itemize}
  \item Which control unit should be chosen for the computer-vision system?
  \item Which camera system should be chosen for computer vision?
  \item Which software solution should be used to solve the problem?
  \item How should information about objects detected by the computer-vision system be passed to the vessel's controllers?
  \item Which methods can be applied to optimise the software?
  \item Which neural network should be used for object detection, and why?
\end{itemize}

\section*{4. Source data}

Source information is gathered from the following sources:

\begin{itemize}
  \item Existing solutions previously built in the world will be studied.
  \item Device data sheets and schematics.
  \item Source information provided by the project team regarding the vessel's control system.
  \item Specialist literature, websites, scientific articles, etc.
\end{itemize}

\section*{5. Research methods}

The theoretical part of the work is based on an analysis of sources (literature, web articles, scientific articles) and on the author's domain knowledge, together with consistency between them.

The practical part covers experiments, comparisons of the models used in computer vision and their capability, taking measurements and analysing them (both graphically and statistically), so that the results obtained enable the most optimal solution to the problem posed.

\section*{6. Graphical part}

Charts, tables, and figures are predominantly in the main body of the work.

\section*{7. Structure of the work}

\textit{(Minor changes to subtopics in the theoretical part are possible.)}

\begin{itemize}
  \item Thesis assignment
  \item Table of contents
  \item Foreword
  \item Glossary of terms, abbreviations, and symbols
  \item Introduction
    \begin{enumerate}
      \item Problem statement
      \item Overview of solutions built around the world so far
      \item Project overview
    \end{enumerate}
  \item Hardware selection, with comparisons
    \begin{enumerate}
      \item Control unit
      \item Cameras
      \item Control-tower assembly drawing / overall diagram
      \item Electrical diagram
    \end{enumerate}
  \item Computer vision and selection of computer-vision software, with comparisons
    \begin{enumerate}
      \item Computer-vision software platforms -- which one, suitable for what
      \item Object detection -- what, how, and why
    \end{enumerate}
  \item Combining computer vision with other sensors for more precise information
    \begin{enumerate}
      \item With radar
      \item With lidar
    \end{enumerate}
  \item Experiments
    \begin{enumerate}
      \item Video-based experiments in a simulated environment to test computer vision
      \item Experiments in the harbour, control-tower test
      \item Experiments on the vessel
      \item Analysis of experimental results
    \end{enumerate}
  \item Conclusions
  \item Summary
  \item Future development work
    \begin{enumerate}
      \item Integration with the AIS (Automatic Identification System)
    \end{enumerate}
  \item List of references
  \item Appendices
\end{itemize}

\section*{8. References used}

Source material consists of specialist literature, maritime-themed articles, and -- with Baltic Workboats' permission -- information on the manoeuvring capabilities of the vessels under test.

\begin{itemize}
  \item Nvidia Jetson AGX Xavier Developer Kit, \url{https://developer.nvidia.com/embedded/jetson-agx-xavier-developer-kit}
  \item Marine Insight, \url{https://www.marineinsight.com/}
  \item OpenCV, \url{https://opencv.org/}
  \item TensorFlow, \url{https://www.tensorflow.org/}
  \item Baltic Workboats, \url{https://bwb.ee/}
  \item MindChip, \url{https://mindchip.ee/}
\end{itemize}

\section*{9. Thesis consultants}

Possible consultation with Baltic Workboats' project manager / client / representative regarding installation of the computer-vision system on the prototype vessel and setting the required input parameters. Also for specifying parameters (turning radius, permitted heel, maximum speed, dimensions from bow to stern and from port to starboard).

\section*{10. Stages and schedule}

\begin{itemize}
  \item Theoretical part (11 March)
    \begin{itemize}
      \item Browsing/collecting sources
      \item Building chapter content
    \end{itemize}
  \item Practical part (2 April)
    \begin{itemize}
      \item Computer-vision model
        \begin{itemize}
          \item Model selection
          \item Implementation
          \item Training
        \end{itemize}
      \item Experiments
        \begin{itemize}
          \item Indoors
          \item Harbour
          \item On the vessel
          \item Developing/refining model parameters
        \end{itemize}
    \end{itemize}
  \item Thesis draft version (23 April)
    \begin{itemize}
      \item With the theoretical part
      \item With the practical part
      \item Submission to the supervisor (24 April)
      \item Making revisions and additions (26 April)
    \end{itemize}
  \item Thesis completed (8 May)
\end{itemize}

\textit{Conditions for a closed defence and/or restrictions on the publication of the thesis are formulated on the reverse.}
